% ECONOMICS_PROBLEM: {"id": "J71-423", "title": "Separating discrimination mechanisms", "jel_code": "J71", "source_id": "OP-0356"}
% !TeX program = xelatex
\documentclass[11pt,a4paper]{article}
\usepackage[margin=23mm]{geometry}
\usepackage{fontspec}
\setmainfont{Latin Modern Roman}
\usepackage{amsmath,amssymb,mathtools}
\usepackage{xcolor,enumitem,booktabs,longtable,array,xurl,hyperref}
\definecolor{muted}{HTML}{53636C}
\hypersetup{colorlinks=true,urlcolor=blue,linkcolor=blue}
\setlength{\parindent}{0pt}
\setlength{\parskip}{5pt}
\newcommand{\R}{\mathbb R}
\newcommand{\E}{\mathbb E}
\newcommand{\N}{\mathbb N}
\newcommand{\ind}{\mathbf 1}
\newcommand{\Var}{\operatorname{Var}}
\newcommand{\Cov}{\operatorname{Cov}}
\newcommand{\supp}{\operatorname{supp}}
\DeclareMathOperator*{\argmin}{arg\,min}
\DeclareMathOperator*{\argmax}{arg\,max}
\newcommand{\entrymeta}[3]{{\sffamily\small\color{muted}\textbf{\texttt{#1}}\quad|\quad #2\par #3\par}}
\newcommand{\Problemref}[1]{\texttt{#1}}
\newcommand{\JELref}[2]{\texttt{#2} (JEL \texttt{#1})}
\begin{document}
\section*{Separating discrimination mechanisms}
\textbf{Problem ID:} \texttt{J71-423}\quad \textbf{JEL:} \texttt{J71}\par
% BEGIN ECONOMICS STATEMENT
\label{problem:OP-0356}
{\sffamily\small\color{muted}Original subject group: Labor markets and work\par}
\label{definitions:problem:30007638}
\textbf{Audit classification:} Background; proposed extension. Authors, title, year and NBER identifier checked. The study separates customer and manager channels in baseball; statistical-discrimination and intersectional variants are editorial.
\subsection*{Definitions and precise research target}
\textbf{Complete editorial problem.} Fix a hiring setting with applicants, employers, and customers. Let \(g_i\) be applicant \(i\)'s recorded identity vector and \(s_i\) independently measured productive characteristics. Let \(H_i(d,c,q)\) indicate hiring under experimentally specified employer information \(d\), customer exposure \(c\), and productivity signal precision \(q\). Employer prejudice concerns identity-dependent preferences holding expected productivity and customer revenue fixed; customer bias concerns the revenue consequences of customer exposure; statistical discrimination concerns decisions based on imperfect productivity beliefs. For identity groups \(g\) and \(g'\), define conditional hiring gaps at common \(s\), and use orthogonal variation in \(d,c,q\) to estimate which channels generate their changes. The identification task is to find assumptions and feasible variations that distinguish these channels, allowing interactions rather than forcing an additive decomposition. A policy arm \(p\) has causal effect \(E[H_i(p)-H_i(0)\mid g_i=g]\), and should also be evaluated for earnings and firm output. Identity labels themselves need not be manipulable. Correspondence tests identify responses to perceived identities under their design; they do not automatically establish motives. Quantify what remains observationally equivalent, especially across intersecting identities and beyond the studied occupation.

\textbf{Definitions and admissible scope.} An identity cue \(\ell\) is a manipulable label that employers perceive; the underlying recorded identity \(g\) is not itself assumed manipulable. Specify employer posterior \(E[s\mid\ell,q]\), customer revenue \(r(s,\ell,c)\), and any identity-dependent employer utility \(v(\ell)\). Distinct triples can generate the same hiring decision, which defines observational equivalence. Let \(H_i(\ell,c,q)\in\{0,1\}\) be the outcome under a complete randomized design. The cue effect \(E[H_i(\ell,c,q)-H_i(\ell',c,q)]\) is identified by randomized cues with common applicant profiles. Separating motives requires restrictions on how \(c\) affects customer demand and how \(q\) affects beliefs; a factorial interaction alone is not a unique channel decomposition. Productivity and revenue measurements must be held comparable across groups. The signal input \(d\) contains the perceived identity cue \(\ell\) and any other randomized applicant information, whereas \(q\) specifies the known precision of a productivity signal and \(c\) the customer-exposure regime. A model supplies the signal experiment, employer prior and update rule, customer-revenue function, employer objective, capacity constraint, and choice rule. At fixed productive profile \(s\), define the recorded-group gap as \(E[H\mid g,s]-E[H\mid g',s]\); this is descriptive unless an admissible intervention and selection assumptions connect it to a causal contrast. The mathematical target is the set of such models generating all observed cell probabilities and policy responses, together with conditions under which the intended channels are separately identified.
\subsection*{Variants and assumption changes}
\begin{enumerate}
\item Customer versus manager channels (source-based scope): restrict to the two channels studied in the cited baseball setting and compare choices across expected customer exposure. Their empirical separation is an existing result, while transport to another industry is an editorial extension.
\item Statistical and intersecting-identity extension (editorial): randomize verified productivity-signal precision and multiple cues jointly. Characterize the set of taste, customer, and belief models fitting all hiring probabilities instead of forcing an additive decomposition.
\end{enumerate}
\subsection*{References and source audit}
\textbf{Support and access.} Authors, title, year and NBER identifier checked. The study separates customer and manager channels in baseball; statistical-discrimination and intersectional variants are editorial. \textbf{Locator:} BFI abstract, January 29, 2025.
\begin{enumerate}\raggedright
\item\hypertarget{definitions:source:30007638:1}{} Majid Ahmadi; Gwen-Jir\={o} Clochard; Jeff Lachman; John A. List. 2025. \emph{Toward an Understanding of Discrimination When Multiple Channels Exist}. Author abstract, opening and final paragraphs, dated January 29, 2025.
\url{https://bfi.uchicago.edu/working-papers/toward-an-understanding-of-discrimination-when-multiple-channels-exist/}.
DOI: \url{https://doi.org/10.3386/w33391}.
\end{enumerate}

\subsection*{Common conventions: Source mathematical conventions}

These conventions apply to each entry unless explicitly replaced.

\textbf{Models and identification.} A primitive model \(m\) specifies all probability laws, feasible choices, information, equilibrium and measurement rules used by its target. The admissible class \(\mathcal M\) is fixed independently of outcomes. If \(P_O(m)\) is its observable probability law and \(\tau(m)\) the target, its identified set at observable law \(P\) is
\[
 \mathcal I_\tau(P)=\{\tau(m):m\in\mathcal M,\ P_O(m)=P\}.
\]
Point identification means that this set is a singleton. An empty set signals incompatibility of data and assumptions. Coordinate bounds need not describe the joint set. Bounds are sharp only if every retained target is attainable or arbitrarily approachable by compatible models. Finite bounds require explicit restrictions on unobserved quantities; unrestricted confounding, hidden wealth, extrapolation or arbitrary equilibrium selection can yield uninformative sets.

\textbf{Causal interventions.} A potential outcome \(Y_i(p)\) is the outcome under a complete policy regime \(p\), including timing, financing, information and market spillovers. The target population and units are fixed before treatment. With interference, use the complete assignment vector or a justified exposure mapping; individual no-interference notation alone cannot identify an economy-wide intervention. A difference of mean potential outcomes does not require the cross-world joint distribution. Conditioning on post-treatment migration, survival, enrollment or employment changes the estimand and needs separate justification.

\textbf{Statistical statements.} An estimator, confidence set, forecast or test specifies a sampling law, model class and loss/coverage criterion. Uniform \((1-\alpha)\)-coverage means \(\inf_{m\in\mathcal M}\Pr_m\{\tau(m)\in C_n\}\geq1-\alpha\), or an explicitly asymptotic version. The whole parameter space trivially has coverage, so informativeness requires a stated width, risk or power objective. A forecasting improvement is relative to a named benchmark, information set and evaluation population.

\textbf{Equilibrium and optimization.} An equilibrium comprises feasible plans, optimality, beliefs where needed, budget constraints and all market/resource clearing. The primitives or a stated benchmark define each correspondence \(\mathcal E(p;m)\). Nonemptiness is assumed or proved, never inferred just from compact policy sets. A selected-equilibrium target uses a declared selection rule; a robust target quantifies over all permitted equilibria. Use suprema or infima unless attainment is established. An \(\varepsilon\)-optimal policy is feasible and within \(\varepsilon\) of the finite optimum value. Report an empty feasible set separately.

\textbf{Welfare and accounting.} Normative utility, interpersonal normalization, social weights, numeraire, discounting and the use of public receipts are primitives. Behavioral utility need not coincide with the welfare criterion. Household welfare includes ownership income and rents once; adding profits again to compensating variation generally double-counts them. A monetary equivalent is defined by an explicit transfer and price convention, with monotonicity and range conditions. Average effects, mechanism contributions, transfers and welfare losses are different targets. Infinite sums require convergence and dynamic feasible plans require terminal or no-Ponzi conditions.

\textbf{Source versions.} A working-paper date, revision date and journal publication year are distinguished. A DOI for a journal version is not automatically the DOI of an earlier manuscript. Locators refer to the stated version and printed pages where specified. A metadata-only check does not verify a claim about a full-text conclusion. Every entry's source subsection narrows the provenance claim accordingly.
% END ECONOMICS STATEMENT
\end{document}
